\section{Khối giao tiếp LoRa và RS485}
Khối giao tiếp trong sơ đồ gồm module LoRa RA-02 và mạch MAX485. Trong phạm vi hệ thống hiện tại, LoRa là kênh truyền chính giữa node và gateway ESP32. MAX485 được xem như kênh mở rộng hoặc dự phòng cho các cảm biến công nghiệp, module ngoại vi hoặc đường truyền có dây trong trường hợp cần độ ổn định cao hơn.

\subsection{Module LoRa RA-02}
RA-02 là module LoRa sử dụng giao tiếp SPI. Các chân quan trọng gồm \texttt{NSS}, \texttt{SCK}, \texttt{MOSI}, \texttt{MISO}, \texttt{RST} và \texttt{DIO0}. Module hoạt động ở mức 3.3V, vì vậy nguồn cấp và mức logic phải được kiểm tra cẩn thận. Trong sơ đồ, các tín hiệu SPI được đưa ra header riêng, giúp quá trình đo và sửa lỗi thuận tiện.

LoRa phù hợp với hệ thống vì dữ liệu cảm biến có dung lượng nhỏ và không yêu cầu băng thông cao. Một gói tin chỉ gồm một số giá trị đo và trạng thái thiết bị, có thể truyền định kỳ sau vài giây hoặc vài chục giây. So với Wi-Fi, LoRa giúp node đơn giản hơn ở lớp mạng, giảm yêu cầu cấu hình và tăng khả năng đặt node ở vị trí xa router.

\subsection{Định dạng gói tin LoRa}
Gói tin cần ngắn gọn, dễ phân tích và có khả năng mở rộng. Định dạng dạng chuỗi phân tách bằng dấu phẩy dễ debug khi dùng Serial Monitor. Tuy nhiên, dạng nhị phân hoặc JSON rút gọn có thể hiệu quả hơn khi cần giảm thời gian truyền. Trong khóa luận, định dạng chuỗi được chọn để dễ kiểm thử.

\begin{lstlisting}[caption={Gói dữ liệu node gửi lên gateway}]
DATA,1,30.9,65.0,2655.0,0.66,698.86,7.1,0,0,0
\end{lstlisting}

Ý nghĩa các trường lần lượt là loại gói, mã node, nhiệt độ, độ ẩm, CO2, ánh sáng, TDS, pH, trạng thái relay 1, relay 2 và điều hòa. Gateway sau khi nhận sẽ chuyển các trường này thành telemetry trên ThingsBoard.

\begin{lstlisting}[caption={Gói lệnh gateway gửi xuống node}]
CMD,1,RL1,ON
CMD,1,RL2,OFF
CMD,1,AC,TEMP_UP
CMD,1,AC,TEMP_DOWN
\end{lstlisting}

Node chỉ thực hiện lệnh khi mã node khớp với địa chỉ của nó và cấu trúc gói tin hợp lệ. Sau khi thực hiện, node gửi lại gói phản hồi:

\begin{lstlisting}[caption={Gói phản hồi trạng thái từ node}]
ACK,1,RL1,ON
ACK,1,AC,TEMP_UP
\end{lstlisting}

\subsection{Cơ chế kiểm tra lỗi}
Trong môi trường căn hộ, LoRa có thể chịu ảnh hưởng bởi vật cản, nhiễu điện từ hoặc khoảng cách. Vì vậy gói tin nên có cơ chế kiểm tra lỗi. Ở mức đơn giản, gateway có thể kiểm tra số trường, kiểu dữ liệu và khoảng giá trị hợp lý. Ở mức tốt hơn, gói tin có thể bổ sung checksum.

\begin{lstlisting}[caption={Ví dụ gói tin có checksum}]
DATA,1,30.9,65.0,2655.0,0.66,698.86,7.1,0,0,0,*5A
\end{lstlisting}

Checksum giúp phát hiện gói tin bị sai trong quá trình truyền. Nếu checksum không khớp, gateway bỏ qua gói tin và không cập nhật telemetry. Đối với lệnh điều khiển, gateway có thể gửi lại nếu không nhận được ACK trong thời gian chờ.

\subsection{Mạch MAX485}
MAX485 là IC chuyển đổi UART sang RS485. RS485 có ưu điểm truyền xa, chống nhiễu tốt và hỗ trợ kết nối nhiều thiết bị trên cùng bus. Trong sơ đồ, MAX485 có các tín hiệu điều khiển hướng truyền/nhận, đường A/B và điện trở kết thúc. Mặc dù hệ thống hiện tại dùng LoRa làm kênh chính, việc bố trí MAX485 giúp node có khả năng mở rộng trong tương lai.

Ví dụ, một cảm biến CO2 công nghiệp hoặc cảm biến môi trường Modbus RTU có thể kết nối qua RS485. Khi đó Arduino Pro Mini sử dụng UART để gửi truy vấn và nhận phản hồi. Tuy nhiên, vì Arduino Pro Mini có tài nguyên hạn chế, cần cân nhắc nếu vừa dùng UART debug, vừa dùng RS485 và vừa dùng LoRa. Trong phiên bản cơ bản, RS485 có thể để dự phòng.

\begin{longtable}{|p{3.2cm}|p{4.5cm}|p{5.5cm}|}
\hline
\textbf{Kênh} & \textbf{Ưu điểm} & \textbf{Hạn chế} \\
\hline
LoRa & Không dây, phạm vi tốt trong căn hộ, phù hợp dữ liệu nhỏ & Băng thông thấp, cần kiểm soát thời gian phát và xác nhận lệnh. \\
\hline
RS485 & Ổn định, chống nhiễu, phù hợp cảm biến công nghiệp & Cần dây dẫn, phức tạp hơn khi lắp đặt trong căn hộ. \\
\hline
Wi-Fi trực tiếp tại node & Kết nối Internet trực tiếp, băng thông cao & Tốn năng lượng hơn, phụ thuộc sóng Wi-Fi tại vị trí đặt node. \\
\hline
\end{longtable}

\subsection{Kiểm thử khối giao tiếp}
Để kiểm thử LoRa, trước hết cần kiểm tra nguồn 3.3V tại module, sau đó kiểm tra SPI bằng chương trình khởi tạo đơn giản. Nếu module không khởi tạo được, cần kiểm tra các chân \texttt{NSS}, \texttt{RST}, \texttt{DIO0}, \texttt{MOSI}, \texttt{MISO} và \texttt{SCK}. Sau khi khởi tạo thành công, thực hiện truyền gói test giữa node và gateway ở khoảng cách ngắn trước, sau đó tăng khoảng cách và thêm vật cản.

Với RS485, kiểm thử gồm đo mức nguồn 5V, kiểm tra tín hiệu DE/RE, kiểm tra đường A/B và thử truyền UART giữa hai thiết bị. Nếu sử dụng điện trở kết thúc, cần bảo đảm giá trị điện trở phù hợp với chiều dài đường dây và cấu hình bus.
